% !TeX program = lualatex
% !TeX encoding = utf8
% !TeX spellcheck = uk_UA
% !TeX root = ../EMConspect.tex

%=========================================================
\chapter{Постійний електричний струм}\label{\currfilebase}
\graphicspath{{\currfilebase/Pictures}}
%=========================================================
\epigraph{Струм --- це не заряд і не поле, а їхня спільна робота: впорядкований рух зарядів, підтримуваний полем, і поле, яке саме
	підтримується перерозподілом зарядів на поверхні провідника.
}{Переосмислення ідей сучасної фізики}


%% --------------------------------------------------------
\section{Основні поняття}
%% --------------------------------------------------------

\subsection*{Сила струму та вектор густини струму}

\emphx{Електричним струмом} називають будь-який упорядкований (спрямований) рух електричних зарядів. \emphx{Силою струму} $I$ називають заряд, що
переноситься через переріз провідника за одиницю часу:
\begin{equation}\label{eq:CurrentDef}
	\tcbhighmath{
		I = \frac{dq}{dt}.
	}
\end{equation}

Струм може бути нерівномірно розподілений по поверхні, крізь яку він протікає (рис.~\ref{pic:IjS}). Щоб описати цей розподіл \emphx{локально}, у
кожній точці простору вводять \emphx{вектор густини струму} $\vect j$: за напрямом він збігається з напрямом руху носіїв \emphx{позитивного} заряду
в даній точці, а за модулем чисельно дорівнює відношенню сили струму через елементарну площадку, розташовану перпендикулярно до руху носіїв, до
площі цієї площадки:
\begin{equation}\label{eq:CurrentDensityDef}
	j = \frac{dI}{dS_\perp}.
\end{equation}

\begin{Attention}
	Напрям вектора $\vect j$ за означенням --- це напрям руху \emphx{позитивних} носіїв заряду. Якщо реальними носіями струму є від'ємно заряджені
	частинки (наприклад, електрони в металі), то вони рухаються \emphx{проти} напряму $\vect j$. Це не є якоюсь фізичною неоднозначністю: напрям
	$\vect j$ --- узгоджена історична домовленість (той самий напрям, що й у "класичного" напряму струму, введеного задовго до відкриття електрона),
	а не твердження про природу носіїв.
\end{Attention}

Знаючи вектор густини струму в кожній точці довільної (не обов'язково перпендикулярної до $\vect j$) поверхні $S$, силу струму через цю поверхню
знаходять як потік вектора $\vect j$:
\begin{equation}\label{eq:CurrentFlux}
	\tcbhighmath{
		I = \iint\limits_S \vect j \cdot d\vect S.
	}
\end{equation}

\begin{figure}[h!]\centering
	\localinput{IjS.tikz}
	\caption{Вектор густини струму $\vect j$ у різних точках поверхні $S$\label{pic:IjS}}
\end{figure}

\subsection*{Густина струму, густина заряду і дрейфова швидкість}

Під дією електричного поля вільні носії заряду в провіднику здійснюють впорядкований рух, який накладається на їхній хаотичний тепловий рух. Середню
швидкість цього впорядкованого руху називають \emphx{дрейфовою швидкістю} $\vect u$.

Нехай концентрація носіїв заряду дорівнює $n$, а заряд кожного носія --- $e$. Візьмемо площадку $dS$, перпендикулярну до дрейфової швидкості
(рис.~\ref{pic:ElectronCylinder}). За час $dt$ цю площадку перетнуть усі носії, що на початку проміжку часу перебували в циліндрі об'ємом
$dV = u\,dt\, dS$: таких носіїв $dN = n\,dV$, і разом вони перенесуть заряд $dq = e\,dN = enu\,dt\,dS$. Отже, густина струму
\begin{equation*}
	j = \frac{dq}{dt}\frac1{dS} = enu = \rho u,
\end{equation*}
де $\rho = en$ --- об'ємна густина заряду носіїв. У векторному вигляді
\begin{equation}\label{eq:JRhoU}
	\tcbhighmath{
		\vect j = \rho\vect u.
	}
\end{equation}

\begin{figure}[h!]\centering
	\localinput{electronCyllinder.tikz}
	\caption{До виведення зв'язку густини струму з дрейфовою швидкістю носіїв\label{pic:ElectronCylinder}}
\end{figure}

\begin{Attention}
	Якщо в провіднику одночасно є носії кількох сортів (наприклад, позитивні й негативні йони в електроліті, або електрони і дірки в
	напівпровіднику), кожен сорт носіїв дає свій внесок у струм, і повна густина струму дорівнює векторній сумі $\vect j = \sum_k \rho_k \vect
	u_k$, причому доданки від носіїв протилежних знаків, які рухаються в протилежних напрямках, \emphx{додаються}, а не віднімаються (адже
	$\rho_k \vect u_k$ не змінює знака при одночасній зміні знаків $\rho_k$ і $\vect u_k$).
\end{Attention}

%% --------------------------------------------------------
\section{Природа носіїв заряду}
%% --------------------------------------------------------

\subsection*{Досліди Толмена і Стюарта}

Той факт, що струм у металах зумовлений рухом \emphx{вільних електронів}, а не, скажімо, йонів кристалічної ґратки, не є очевидним апріорі і
потребував експериментального доведення. У 1916~р. американський фізик Р.~Толмен і шотландський фізик Т.~Стюарт виконали дослід, який
неспростовно це довів.

Котушку з великою кількістю витків тонкого дроту приєднували до гальванометра і приводили в швидке обертання навколо власної осі, а потім різко
гальмували. Під час гальмування вільні носії заряду, за інерцією, продовжували свій рух відносно ґратки провідника ще якийсь час, тобто в колі
виникав короткочасний струм. За напрямом відхилення стрілки гальванометра було встановлено, що цей струм створюють \emphx{негативно заряджені}
частинки, а виміряний питомий заряд носіїв $e/m$ виявився близьким до питомого заряду електрона, відомого з інших дослідів (зокрема, з дослідів Дж.
Дж. Томсона з катодними променями). Так було експериментально доведено, що носіями струму в металах є електрони.

\begin{Attention}
	Ідея досліду Толмена і Стюарта загалом однакова з ідеєю досліду з маятником Фуко чи гіроскопом: інерція носіїв заряду --- це прояв того самого
	другого закону Ньютона, застосованого не до макроскопічного тіла, а до \emphx{мікроскопічних} носіїв заряду всередині нього. Такий підхід,
	коли макроскопічно вимірювана величина (тут --- імпульсний струм) дає інформацію про мікроскопічну характеристику речовини (тут --- питомий
	заряд носіїв), типовий для фізики початку XX~століття.
\end{Attention}

\subsection*{Приклад: дрейфова і теплова швидкості електронів}

Порівняємо на конкретному прикладі дрейфову швидкість електронів (швидкість їхнього впорядкованого руху під дією поля) з їхньою тепловою швидкістю
(швидкістю хаотичного теплового руху) --- порівняння, яке добре ілюструє, наскільки "повільним" насправді є струм на мікроскопічному рівні.

Нехай срібним дротом перерізом $S=1$~мм$^2$ тече струм силою $I=1$~А; вважатимемо, що кожен атом срібла віддає один вільний електрон. Струм
$I = jS = enuS$, а концентрацію електронів знаходимо через густину срібла $\rho_m$, його молярну масу $\mu$ і сталу Авогадро $N_A$:
\begin{equation*}
	n = \frac{\rho_m}{m_{\ce{Ag}}} = \rho_m\frac{N_A}{\mu}.
\end{equation*}
Звідси дрейфова швидкість
\begin{equation*}
	u = \frac{I}{enS} = \frac{I\mu}{e\rho_m N_A S} \approx 0.065\ \text{мм/с}.
\end{equation*}
Теплову ж швидкість електронів оцінимо як середню квадратичну швидкість молекул ідеального газу при кімнатній температурі $T=300$~K:
\begin{equation*}
	v_T = \sqrt{\frac{3kT}{m_e}} \approx 1.2\cdot 10^5\ \text{м/с}.
\end{equation*}

\begin{Attention}
	Дрейфова швидкість електронів у типовому провіднику --- частки міліметра за секунду, тобто на \emphx{дев'ять порядків} менша за теплову
	швидкість електронів! Струм у колі "вмикається" практично миттєво (з електромагнітною хвилею, зі швидкістю порядку швидкості світла) не тому,
	що електрони "біжать" зі старту до навантаження, а тому, що електричне поле, яке приводить в рух \emphx{усі} вільні електрони провідника
	\emphx{одночасно} (а не послідовно, один за одним), встановлюється в колі майже миттєво.
\end{Attention}

%% --------------------------------------------------------
\section{Закон збереження заряду}
%% --------------------------------------------------------

Нехай у деякому провідному середовищі виділено об'єм $V$, обмежений замкненою поверхнею $S$. Повний струм, що витікає крізь цю поверхню,
$I = \oiint_S \vect j\cdot d\vect S$; оскільки струм --- це рух заряду, то заряд, який лишається всередині об'єму $V$, з часом \emphx{зменшується}
рівно на ту величину заряду, яка витекла (рис.~\ref{pic:ChargeConservation}):
\begin{equation}\label{eq:ContinuityInt}
	\oiint\limits_S \vect j\cdot d\vect S = -\frac{d}{dt}\iiint\limits_V \rho\, dV.
\end{equation}
Це співвідношення називають \emphx{рівнянням неперервності}: воно є математичним виразом закону збереження електричного заряду (заряд не
з'являється й не зникає --- він лише переноситься з місця на місце). У диференціальній формі (за теоремою Остроградського--Гаусса)
\begin{equation}\label{eq:ContinuityDiff}
	\tcbhighmath{
		\divg \vect j = -\frac{\partial\rho}{\partial t}.
	}
\end{equation}

\begin{figure}[h!]\centering
	\localinput{chargeConservation.tikz}
	\caption{Струм, що витікає крізь замкнену поверхню $S$, зменшує заряд в об'ємі $V$\label{pic:ChargeConservation}}
\end{figure}

У \emphx{стаціонарному} випадку, тобто коли розподіл заряду в просторі не змінюється з часом ($\partial\rho/\partial t = 0$), рівняння неперервності
спрощується:
\begin{equation}\label{eq:ContinuitySteady}
	\oiint\limits_S \vect j\cdot d\vect S = 0, \qquad \text{або}\qquad \divg\vect j = 0.
\end{equation}
Це означає, що з будь-якого об'єму, обмеженого замкненою поверхнею, витікає рівно стільки ж заряду, скільки в нього і втікає --- лінії струму в
стаціонарному випадку не мають ні витоків, ні стоків (окрім хіба що точок, де є зовнішні джерела чи стоки заряду).

\begin{Attention}
	Саме умова $\divg \vect j = 0$ --- це те, що математично означає слова "\emphx{постійний} (стаціонарний) струм": лінії струму в такому режимі
	замкнені (або йдуть з нескінченності в нескінченність), і сила струму однакова в будь-якому перерізі \emphx{послідовно} з'єднаного провідника
	--- саме тому має сенс говорити про "силу струму в колі", не уточнюючи, в якому саме місці кола вона виміряна.
\end{Attention}

%% --------------------------------------------------------
\section{Закон Ома}
%% --------------------------------------------------------

\subsection*{Інтегральна форма}

Г.~Ом 1827~року експериментально встановив, що сила струму $I$, який тече однорідним металевим провідником, за відсутності в ньому сторонніх сил,
пропорційна різниці потенціалів (напрузі) $U$ на його кінцях:
\begin{equation}\label{eq:OhmIntegral}
	\tcbhighmath{
		I = \frac{U}{R}.
	}
\end{equation}
Коефіцієнт пропорційності $R$ називають \emphx{електричним опором} провідника; він залежить від форми й розмірів провідника, від матеріалу, з якого
він виготовлений, від температури, а також, узагалі кажучи, від розподілу струму по провіднику. У найпростішому випадку однорідного циліндричного
провідника довжиною $\ell$ і перерізом $S$
\begin{equation}\label{eq:ResistanceCylinder}
	R = \rho\, \frac{\ell}{S},
\end{equation}
де $\rho$ --- \emphx{питомий опір} матеріалу (величина, обернена до питомої електропровідності $\lambda = 1/\rho$, яку ми означимо нижче).

\subsection*{Про одиниці вимірювання}

У гауссовій системі одиниць $[I] = \text{Фр}/\text{с}$, а $[U] = \text{Фр}/\text{см}$, тому
\begin{equation*}
	[R] = \frac{[U]}{[I]} = \frac{\text{с}}{\text{см}},
\end{equation*}
тобто опір у СГС має розмірність \emphx{оберненої швидкості}. Відповідно, для питомого опору
\begin{equation*}
	[\rho] = [R]\frac{[S]}{[\ell]} = \frac{\text{с}}{\text{см}}\cdot\frac{\text{см}^2}{\text{см}} = \text{с},
\end{equation*}
тобто в гауссовій системі питомий опір вимірюється просто в \emphx{секундах}, а питома провідність $\lambda$ --- в обернених секундах.

\begin{Attention}
	Незвична на перший погляд розмірність опору й питомого опору в гауссовій системі одиниць --- не якийсь дефект цієї системи, а наслідок того, що
	в СГС немає окремої "електричної" основної одиниці (на відміну від Ампера в СІ): усі електричні величини виражаються через звичні механічні
	одиниці --- сантиметри, грами, секунди. У Міжнародній системі одиниць (СІ) опір, звісно, вимірюють в Омах, а питомий опір --- в Ом$\cdot$м; саме
	ці одиниці й наведено в таблиці питомих опорів металів нижче.
\end{Attention}

\subsection*{Диференціальна форма}

Виведемо \emphx{локальний} (диференціальний) варіант закону Ома, який пов'язує густину струму $\vect j$ з напруженістю поля $\Efield$ \emphx{в
	кожній точці} провідника, а не інтегрально --- на його кінцях. Виділимо в околі довільної точки провідного середовища нескінченно малий
циліндричний об'єм із твірними, паралельними вектору $\vect j$ (рис.~\ref{pic:DiffOhmLaw}).

\begin{figure}[h!]\centering
	\localinput{diffOhmLow.tikz}
	\caption{Елементарний циліндр струму до виведення диференціальної форми закону Ома\label{pic:DiffOhmLaw}}
\end{figure}

Якщо переріз циліндра дорівнює $dS$, а довжина --- $dl$, то, застосовуючи до нього закон
Ома~\eqref{eq:OhmIntegral}--\eqref{eq:ResistanceCylinder} ($I \to j\, dS$, $U \to E\, dl$, $R \to \rho\, dl/dS$), маємо
\begin{equation*}
	j\, dS = \frac{E\, dl}{\rho\, dl/dS} = \frac1\rho E\, dS,
\end{equation*}
звідки, після скорочень, у векторному вигляді:
\begin{equation}\label{eq:OhmDiff}
	\tcbhighmath{
		\vect j = \frac1\rho \Efield = \lambda\Efield,
	}
\end{equation}
де $\lambda = 1/\rho$ --- питома електропровідність середовища.

\subsection*{Теорія провідності металів Друде}

У 1900~р. П.~Друде запропонував просту класичну модель, що якісно (а часто й кількісно правильно за порядком величини) пояснює електропровідність
металів. У цій моделі вільні електрони металу розглядають як класичний "газ", що рухається в полі $\Efield$ і час від часу зазнає зіткнень з
іонами кристалічної ґратки, розсіюючись на них.

Рівняння руху окремого електрона між зіткненнями: $\vect a = e\Efield/m_e$. Якщо одразу після зіткнення швидкість електрона (в середньому)
дорівнює нулю, то до наступного зіткнення вона зростає лінійно з часом, $v = at$ (рис.~\ref{pic:Ut}); після зіткнення набута впорядкована швидкість,
в середньому, знову "обнуляється", і цикл прискорення повторюється.

\begin{figure}[h!]\centering
	\localinput{ut.tikz}
	\caption{Пилкоподібна залежність швидкості електрона від часу в моделі Друде\label{pic:Ut}}
\end{figure}

Якщо середній час вільного пробігу дорівнює $\tau$, то середня (за час між зіткненнями) швидкість упорядкованого руху
\begin{equation*}
	u = \overline v = \frac12\, \frac{e E \tau}{m_e},
\end{equation*}
звідки, за формулою~\eqref{eq:JRhoU}, $j = enu = \dfrac{ne^2\tau}{2m_e} E$, тобто питома провідність за теорією Друде
\begin{equation}\label{eq:DrudeConductivity}
	\tcbhighmath{
		\lambda = \frac{ne^2\tau}{2m_e}.
	}
\end{equation}

\begin{Attention}
	Множник $\tfrac12$ у формулі~\eqref{eq:DrudeConductivity} --- наслідок \emphx{спрощеного} припущення, що всі електрони пролітають між
	зіткненнями рівно однаковий час $\tau$ (звідси й "пилкоподібний" графік $v(t)$, лінійний на кожній ділянці, із середнім значенням рівно
	посередині). Насправді ж час вільного пробігу --- випадкова величина з (приблизно) експоненційним розподілом, і акуратне усереднення за цим
	розподілом \emphx{прибирає} множник $1/2$, даючи $\lambda = ne^2\tau/m_e$, де $\tau$ тепер --- \emphx{середній} час вільного пробігу.
	Розбіжність рівно вдвічі --- класична "пастка" елементарного виведення теорії Друде, на яку варто звертати увагу: за порядком величини обидві
	формули однаково добре узгоджуються з експериментом (сама теорія Друде --- лише грубе класичне наближення, оскільки насправді електрони в
	металі підпорядковуються квантовій статистиці Фермі--Дірака), тому конкретний числовий множник тут не варто сприймати надто серйозно.
\end{Attention}

\subsection*{Температурна залежність опору}

Час вільного пробігу $\tau$, а отже, і провідність $\lambda$, залежать від температури, оскільки з підвищенням температури зростає інтенсивність
теплових коливань іонів ґратки, на яких розсіюються електрони. Найпростіший варіант теорії Друде (у припущенні, що концентрація $n$ від температури
не залежить, а частота зіткнень пропорційна швидкості теплового руху електронів $v_T \propto \sqrt T$) передбачає
\begin{equation*}
	\rho \propto \sqrt T
\end{equation*}
(рис.~\ref{pic:DrudeRho}, зліва). Реальна залежність питомого опору чистих металів від температури для не надто низьких температур --- майже
\emphx{лінійна} (рис.~\ref{pic:ExpRho}, справа):
\begin{equation}\label{eq:TempDependence1}
	\rho = \rho_0\left(1 + \alpha T\right)
\end{equation}
або, якщо відлічувати температуру від зручного початку відліку $20\,^\circ$C,
\begin{equation}\label{eq:TempDependence2}
	\rho = \rho_{20}\left[1 + \alpha(t - 20)\right].
\end{equation}

\begin{figure}[h!]
	\centering
	\begin{minipage}[t]{0.45\linewidth}\centering
		\localinput{DrudeRho.tikz}
		\caption{Передбачення найпростішої теорії Друде: $\rho\propto\sqrt T$}\label{pic:DrudeRho}
	\end{minipage}\hfill
	\begin{minipage}[t]{0.45\linewidth}\centering
		\localinput{ExpRho.tikz}
		\caption{Типова експериментальна залежність: лінійна ділянка й стрибок у нуль при
			надпровідному переході}\label{pic:ExpRho}
	\end{minipage}
\end{figure}

\begin{table}[h!]\small\centering
	\caption{Питомий опір деяких металів при $20\,^\circ$C}
	\begin{tabular}{lcc}
		\toprule
		Метал              & $\rho,\ 10^{-8}\,\text{Ом}\cdot\text{м}$ & $\alpha,\ 10^{-3}\,{}^\circ\text{C}^{-1}$ \\
		\midrule
		Мідь (\ce{Cu})     & 1.68                                      & 4.3                                        \\
		Алюміній (\ce{Al}) & 2.82                                      & 4.2                                        \\
		Залізо (\ce{Fe})   & 9.71                                      & 6.0                                        \\
		Нікель (\ce{Ni})   & 6.84                                      & 6.5                                        \\
		Вольфрам (\ce{W})  & 5.65                                      & 5.0                                        \\
		\bottomrule
	\end{tabular}
\end{table}

\begin{Attention}
	Лінійне зростання опору з температурою характерне саме для \emphx{металів}, у яких концентрація вільних носіїв $n$ практично не залежить від
	температури, а опір визначається тільки частотою зіткнень. У \emphx{напівпровідниках}, навпаки, з ростом температури різко (експоненційно)
	зростає концентрація вільних носіїв, тому опір напівпровідників, як правило, з ростом температури \emphx{спадає} --- поведінка, протилежна до
	металів. Ще одне яскраве явище, показане на рис.~\ref{pic:ExpRho} стрибком до нуля, --- \emphx{надпровідність}: у деяких матеріалах при
	охолодженні нижче деякої критичної температури $T_c$ опір стрибкоподібно (а не поступово) звертається в точний нуль.
\end{Attention}

\subsection*{Розподіл зарядів у провіднику, яким тече струм}

З'ясуємо, де саме розташовані заряди, які створюють електричне поле всередині провідника зі струмом. Якщо провідник \emphx{однорідний}
($\lambda = \const$) і по ньому тече стаціонарний струм, то $\divg\vect j = 0$ за формулою~\eqref{eq:ContinuitySteady}. Але, за диференціальним
законом Ома~\eqref{eq:OhmDiff} і теоремою Гаусса,
\begin{equation}\label{eq:VolumeChargeZero}
	\divg\vect j = \lambda\,\divg\Efield = \lambda\cdot 4\pi\rho = 0 \quad\Rightarrow\quad \rho = 0.
\end{equation}
Отже, в однорідному провіднику зі стаціонарним струмом \emphx{об'ємна} густина заряду в усіх внутрішніх точках дорівнює нулю.

Водночас, на \emphx{поверхні} провідника зі струмом заряди, взагалі кажучи, є (рис.~\ref{pic:SurfCharges}). Саме ці поверхневі заряди і є джерелами
поля $\Efield$, яке існує всередині провідника і забезпечує протікання постійного струму: усередині ідеального провідника без струму поле дорівнює
нулю (як ми вже обговорювали в розділі про провідники в електростатичній рівновазі), але для \emphx{підтримання} струму поле всередині мусить бути відмінним від нуля, і саме поверхневі заряди
перерозподіляються так, щоб створити потрібну конфігурацію поля.

\begin{figure}[h!]\centering
	\localinput{SurfCharges.tikz}
	\caption{Поверхневі заряди на провіднику зі струмом створюють як тангенціальну, так і нормальну
		складову поля\label{pic:SurfCharges}}
\end{figure}

\begin{Attention}
	Це принципова відмінність провідника зі струмом від провідника в електростатиці: в електростатиці поверхня провідника завжди
	\emphx{еквіпотенціальна}, а поле всередині відсутнє. У провіднику зі струмом поверхневий заряд перерозподіляється так, щоб забезпечити не
	відсутність поля всередині, а якраз \emphx{потрібне для підтримання струму} поле $\Efield = \vect j/\lambda$ у кожній внутрішній точці --- і
	тому поверхня провідника зі струмом, узагалі кажучи, вже \emphx{не} еквіпотенціальна.
\end{Attention}

\subsection*{Сторонні сили. Електрорушійна сила}

Щоб у колі підтримувався постійний струм, недостатньо ділянок, на яких позитивні носії рухаються в бік \emphx{спадання} потенціалу (тобто під дією
електричного поля): мають бути й ділянки, де носії рухаються в бік \emphx{зростання} потенціалу, тобто \emphx{проти} сил електростатичного поля
(рис.~\ref{pic:CurrentSpheres}). Перенесення носіїв на таких ділянках можливе лише за допомогою сил \emphx{неелектростатичного} походження, які
називають \emphx{сторонніми силами}.

\begin{figure}[h!]\centering
	\localinput{CurrentSpheres.tikz}
	\caption{У замкненому колі частина носіїв мусить рухатися проти сил електростатичного поля --- це
		можливо лише завдяки стороннім силам\label{pic:CurrentSpheres}}
\end{figure}

Фізична природа сторонніх сил може бути дуже різною: хімічна енергія (гальванічні елементи, акумулятори), різниця температур (термоелементи),
енергія світла (фотоелементи), механічна енергія (генератори) тощо. Незалежно від конкретної природи, для кількісного опису вводять поле сторонніх
сил $\vect E^*$ --- сторонню силу, що діє на одиничний позитивний заряд:
\begin{equation}\label{eq:ExtraneousForceField}
	\vect E^* = \frac{\vect F^*}{q}.
\end{equation}
Роботу сторонніх сил з переміщення одиничного позитивного заряду з точки $1$ у точку $2$ (в межах ділянки, де діють сторонні сили) називають
\emphx{електрорушійною силою} (ЕРС):
\begin{equation}\label{eq:EMFDef}
	\tcbhighmath{
		\mathcal E = \frac{A^*}{q} = \int\limits_{(1)}^{(2)} \vect E^* \cdot d\vect r.
	}
\end{equation}

\begin{Attention}
	Незважаючи на назву, ЕРС --- це не сила, а робота (точніше, робота на одиницю заряду), тобто величина з розмірністю потенціалу (напруги), а не
	сили. Історична назва --- данина часів, коли поняття енергії й роботи в електриці ще не було чітко відокремлено від поняття сили.
\end{Attention}

Якщо на провідник одночасно діють і електростатичне поле $\Efield$, і поле сторонніх сил $\vect E^*$, то на кожен носій заряду діє їхня сума, і
диференціальний закон Ома~\eqref{eq:OhmDiff} узагальнюється:
\begin{equation}\label{eq:GeneralizedOhm}
	\tcbhighmath{
		\vect j = \lambda\left(\Efield + \vect E^*\right).
	}
\end{equation}

\subsection*{Закон Ома в інтегральній формі для неоднорідної ділянки і замкненого кола}

\emphx{Неоднорідною} називають ділянку кола, на якій, окрім електростатичного поля, діють ще й сторонні сили (рис.~\ref{pic:OhmLaw}). Розглянемо
таку ділянку між точками $1$ і $2$ і проінтегруємо узагальнений закон Ома~\eqref{eq:GeneralizedOhm}, скалярно помножений на елемент довжини $d\vect
r$ уздовж ділянки, від точки $1$ до точки $2$:
\begin{equation*}
	\int\limits_{(1)}^{(2)} \frac{\vect j\cdot d\vect r}{\lambda} = \int\limits_{(1)}^{(2)} \Efield\cdot d\vect r + \int\limits_{(1)}^{(2)} \vect
	E^*\cdot d\vect r.
\end{equation*}

\begin{figure}[h!]\centering
	\localinput{OhmLow.tikz}
	\caption{Неоднорідна ділянка кола: точка $1$ (потенціал $\phi_1$) і точка $2$ (потенціал
		$\phi_2$), між якими діють і опір, і сторонні сили\label{pic:OhmLaw}}
\end{figure}

Ліва частина, з огляду на сталість сили струму $I$ уздовж послідовно з'єднаної ділянки (стаціонарність!) і означення опору, дорівнює $IR$, де
$R = \int \rho\, dl/S$ --- опір усієї ділянки; перший доданок справа, за означенням потенціалу, дорівнює $\phi_1 - \phi_2$; другий, за
означенням~\eqref{eq:EMFDef}, --- це ЕРС $\mathcal E$, що діє на ділянці. Отримуємо \emphx{закон Ома для неоднорідної ділянки кола}:
\begin{equation}\label{eq:OhmNonUniform}
	\tcbhighmath{
		IR = \phi_1 - \phi_2 + \mathcal E.
	}
\end{equation}

\begin{Attention}
	Знак ЕРС у формулі~\eqref{eq:OhmNonUniform} --- \emphx{алгебраїчний}: якщо джерело ЕРС на цій ділянці сприяє руху позитивних носіїв в обраному
	напрямку (від $1$ до $2$), то $\mathcal E > 0$; якщо перешкоджає --- $\mathcal E < 0$. Практичне правило: обходячи ділянку від $1$ до $2$, ЕРС
	беруть зі знаком "плюс", якщо цей обхід усередині джерела відбувається від мінуса до плюса (тобто в напрямку дії сторонніх сил), і зі знаком
	"мінус" --- у протилежному випадку.
\end{Attention}

Для \emphx{замкненого} кола точки $1$ і $2$ збігаються, тому $\phi_1 = \phi_2$, і закон~\eqref{eq:OhmNonUniform} набуває найпростішого вигляду:
\begin{equation}\label{eq:OhmClosedLoop}
	\tcbhighmath{
		IR = \mathcal E,
	}
\end{equation}
де тепер $R$ --- повний опір усього замкненого кола, а $\mathcal E$ --- алгебраїчна сума всіх ЕРС, що діють у цьому колі.

\subsection*{Розгалужені кола. Правила Кірхгофа}

Для розрахунку розгалужених кіл (тобто кіл з кількома контурами, що мають спільні ділянки) закони~\eqref{eq:OhmNonUniform}--\eqref{eq:OhmClosedLoop}
зручно записувати у формі двох правил Кірхгофа.

\emphx{Перше правило Кірхгофа} стосується \emphx{вузлів} кола --- точок, де сходяться три й більше провідники (рис.~\ref{pic:Nodes}) --- і є прямим
наслідком стаціонарності струму~\eqref{eq:ContinuitySteady}, застосованого до нескінченно малого об'єму навколо вузла: алгебраїчна сума струмів, що
сходяться у вузлі, дорівнює нулю:
\begin{equation}\label{eq:Kirchhoff1}
	\tcbhighmath{
		\sum\limits_{k=1}^K I_k = 0.
	}
\end{equation}
Струми, що втікають у вузол, і струми, що витікають з нього, при цьому слід вважати величинами \emphx{різних знаків} (наприклад, "втікає" --- зі
знаком "плюс", "витікає" --- зі знаком "мінус", або навпаки, аби послідовно дотримуватись однієї домовленості).

\begin{figure}[h!]\centering
	\localinput{nodes.tikz}
	\caption{Вузол розгалуженого кола: $I_1 + I_4 = I_2 + I_3$\label{pic:Nodes}}
\end{figure}

\emphx{Друге правило Кірхгофа} стосується довільного \emphx{контуру} (замкненої частини розгалуженого кола, рис.~\ref{pic:Loop}) і є прямим
наслідком закону Ома для замкненого кола~\eqref{eq:OhmClosedLoop}, застосованого до цього контуру: алгебраїчна сума добутків сил струмів на окремих
ділянках контуру на їхні опори дорівнює алгебраїчній сумі ЕРС, що діють у цьому контурі:
\begin{equation}\label{eq:Kirchhoff2}
	\tcbhighmath{
		\sum\limits_{m=1}^M I_m R_m = \sum\limits_{n=1}^N \mathcal E_n.
	}
\end{equation}

\begin{figure}[h!]\centering
	\localinput{loop.tikz}
	\caption{Контур розгалуженого кола з трьома ЕРС і трьома резисторами\label{pic:Loop}}
\end{figure}

\begin{Attention}
	Щоб послідовно застосувати друге правило Кірхгофа, потрібно: (1)~довільно обрати напрям обходу контуру (за годинниковою стрілкою чи проти);
	(2)~довільно обрати (і позначити на схемі) напрями невідомих струмів $I_m$ на кожній ділянці; (3)~доданок $I_mR_m$ брати зі знаком "плюс", якщо
	обраний напрям струму на цій ділянці збігається з напрямом обходу, і зі знаком "мінус" --- якщо ні; (4)~ЕРС $\mathcal E_n$ брати зі знаком
	"плюс", якщо напрям обходу контуру всередині джерела йде від мінуса до плюса, і зі знаком "мінус" --- у протилежному випадку. Якщо в результаті
	розв'язання системи рівнянь якийсь із струмів $I_m$ виявиться від'ємним, це означає лише, що його реальний напрям --- протилежний до обраного
	на кроці (2), а не помилку в розрахунку.
\end{Attention}

%% --------------------------------------------------------
\section{Закон Джоуля--Ленца}
%% --------------------------------------------------------

Оскільки провідник зі струмом чинить опір рухові носіїв заряду (розсіюючи їхню впорядковану кінетичну енергію на теплові коливання ґратки при
зіткненнях), проходження струму неминуче супроводжується виділенням тепла.

\emphx{Мікроскопічне виведення (за теорією Друде).} При зіткненні електрона з іоном уся кінетична енергія, набута електроном за час вільного
пробігу, $\tfrac12 m_e v_{\max}^2$ (де $v_{\max} = eE\tau/m_e$ --- швидкість електрона безпосередньо перед зіткненням), передається ґратці у вигляді
тепла. Число зіткнень одного електрона за одиницю часу дорівнює $1/\tau$, тому число зіткнень за одиницю часу в одиниці об'єму дорівнює $N = n/\tau$
($n$ --- концентрація електронів), а кількість теплоти, що виділяється в одиниці об'єму провідника за одиницю часу,
\begin{equation*}
	w = N\cdot\frac{m_e v_{\max}^2}{2} = \frac n\tau\cdot\frac{m_e}2\left(\frac{eE\tau}{m_e}\right)^2 = \frac{ne^2\tau}{2m_e}E^2 = \lambda E^2.
\end{equation*}

\emphx{Загальне (енергетичне) виведення.} Той самий результат можна отримати значно простіше й зовсім не апелюючи до мікроскопічної моделі. Робота
поля з перенесення заряду $dq$ через різницю потенціалів $U$ дорівнює $dA = U\, dq = UI\, dt$; уся ця робота, за відсутності інших змін в системі
(наприклад, накопичення енергії в конденсаторі), переходить у теплоту, тому потужність тепловиділення на ділянці кола $\mathcal P = dA/dt = UI$, а
за законом Ома $U=IR$:
\begin{equation*}
	\mathcal P = I^2 R = \frac{U^2}{R}.
\end{equation*}
Записуючи це для нескінченно малого циліндричного елемента провідника ($U \to E\, dl$, $I \to j\, dS$) і ділячи на його об'єм $dV = dl\, dS$,
отримуємо ту саму формулу для \emphx{об'ємної густини потужності тепловиділення}:
\begin{equation}\label{eq:JouleLenzDiff}
	\tcbhighmath{
		w = \lambda E^2 = \vect j\cdot\Efield.
	}
\end{equation}
Інтегруючи по всьому об'єму провідника, знаходимо \emphx{закон Джоуля--Ленца в інтегральній формі}:
\begin{equation}\label{eq:JouleLenzIntegral}
	\tcbhighmath{
		\mathcal P = \int\limits_V w\, dV = \int\limits_V \rho j^2\, dV = I^2\int\limits_V \rho\,\frac{dl}{S} = I^2 R.
	}
\end{equation}

\begin{Attention}
	Той факт, що обидва підходи --- детальний мікроскопічний (через зіткнення окремих електронів) і гранично загальний макроскопічний (через
	закон збереження енергії) --- дають \emphx{однаковий} результат, зовсім не випадковий: закон Джоуля--Ленца є прямим наслідком закону збереження
	енергії і тому не залежить від деталей мікроскопічного механізму розсіювання --- він буде справедливим і для електронної провідності в металах,
	і для йонної провідності в електролітах, і для будь-якого іншого механізму опору.
\end{Attention}

%% --------------------------------------------------------
\section{Струми в необмежених середовищах}
%% --------------------------------------------------------

Розглянемо задачу, яка на перший погляд не стосується "кіл" у звичному розумінні: у провідне середовище з питомою провідністю $\lambda$ і
діелектричною проникністю $\epsilon$ занурено два електроди, і потрібно знайти повний опір середовища між ними (рис.~\ref{pic:MediaCurrent}). Ця
задача важлива, наприклад, для розрахунку опору заземлення.

\begin{figure}[h!]\centering
	\localinput{MediaCurrent.tikz}
	\caption{Два електроди в необмеженому провідному середовищі з провідністю $\lambda$ і
		проникністю $\epsilon$\label{pic:MediaCurrent}}
\end{figure}

Нехай на електроди нанесено заряди $+q$ і $-q$. З теореми Гаусса для електростатичного поля в діелектрику, застосованої до поверхні, що охоплює
позитивний електрод,
\begin{equation*}
	\oiint\limits_S \Efield\cdot d\vect S = \frac{4\pi}{\epsilon}\, q = \frac{4\pi}{\epsilon}\, C(\phi_+ - \phi_-),
\end{equation*}
де $C$ --- взаємна ємність електродів (у середовищі з проникністю $\epsilon$), а $q = C(\phi_+ - \phi_-)$. Але та сама поверхня охоплює й струм, що
стікає з електрода: за диференціальним законом Ома, $\vect j = \lambda\Efield$ поблизу поверхні електрода, тому повний струм, що стікає з нього,
\begin{equation*}
	I = \oiint\limits_S \vect j\cdot d\vect S = \lambda\oiint\limits_S \Efield\cdot d\vect S = \frac{4\pi\lambda}{\epsilon}\, C(\phi_+ - \phi_-).
\end{equation*}
Звідси опір середовища між електродами:
\begin{equation}\label{eq:GroundResistanceGeneral}
	\tcbhighmath{
		R = \frac{\phi_+ - \phi_-}{I} = \frac{\epsilon}{4\pi\lambda C}.
	}
\end{equation}

\begin{Attention}
	Формула~\eqref{eq:GroundResistanceGeneral} виражає чудову й зовсім не очевидну \emphx{аналогію} між електростатичною задачею про ємність і
	задачею про опір провідного середовища: обидві задачі зводяться до розв'язання \emphx{однакового} рівняння (рівняння Лапласа для потенціалу)
	з однаковими граничними умовами на поверхнях електродів --- відрізняється лише те, яку саме величину ми обчислюємо через потенціал ($q$ чи
	$I$). Тому \emphx{щойно відома ємність} якоїсь конфігурації електродів, опір середовища між ними отримують одразу, без жодних додаткових
	обчислень, за формулою~\eqref{eq:GroundResistanceGeneral}.
\end{Attention}

Скориставшись цією аналогією, відразу випишемо два важливі окремі випадки.

\emphx{Дві далекі кулі-електроди.} Якщо електроди --- провідні кулі з власними ємностями $C_1 = \epsilon r_1$ і $C_2 = \epsilon r_2$, віддалені
одна від одної на велику відстань, то, за формулою для взаємної ємності далеких провідників (гл.~\ref{ElEnergy}),
$1/C = 1/C_1 + 1/C_2 = (1/\epsilon)(1/r_1+1/r_2)$, тому
\begin{equation}\label{eq:GroundResistanceTwoElectrodes}
	R = \frac1{4\pi\lambda}\left(\frac1{r_1}+\frac1{r_2}\right),
\end{equation}
--- тобто опір заземлення практично \emphx{не залежить} від відстані між кулями (вона взагалі не входить у формулу), а визначається лише розмірами
самих електродів.

\emphx{Плоскі електроди.} Якщо електроди --- обкладки плоского конденсатора площею $S$ на відстані $d$ одна від одної, то $C = \epsilon S/(4\pi d)$,
і
\begin{equation}\label{eq:GroundResistancePlates}
	R = \frac{d}{\lambda S},
\end{equation}
що просто збігається з відомою формулою~\eqref{eq:ResistanceCylinder} для опору однорідного провідника довжиною $d$, перерізом $S$ і питомим опором
$1/\lambda$ --- як, зрештою, і мало бути.

\subsection*{Приклад: опір заземлення і крокова напруга}

Нехай фундамент металевої опори лінії електропередач має вигляд півсфери діаметром $D=2$~м, а ґрунт навколо неї має провідність
$\lambda = 2\cdot10^{-4}$~См/см і слугує заземленням (рис.~\ref{pic:StepVoltage}). Знайдемо опір заземлення та \emphx{крокову напругу} на відстані
$r=5$~м від центру опори в разі замикання проводу високої напруги $\phi_0=10$~кВ на опору, якщо довжина кроку людини $\ell=0.7$~м.

\begin{figure}[h!]\centering
	\localinput{StepVoltage.tikz}
	\caption{Розтікання струму від напівсферичного заземлювача та крокова напруга\label{pic:StepVoltage}}
\end{figure}

Ємність напівсфери радіуса $D/2$, що межує із землею, вдвічі більша за ємність повної кулі того самого радіуса в необмеженому середовищі (оскільки
поле існує лише в півпросторі), тому $C = \epsilon D$ (у гауссовій системі, для землі $\epsilon\approx 1$), і за
формулою~\eqref{eq:GroundResistanceGeneral} з одним електродом на нескінченності ($C_2\to\infty$, тобто $1/C \to 1/C_1$)
\begin{equation*}
	R = \frac{1}{\pi\lambda D} \approx 8\ \text{Ом}.
\end{equation*}
Потенціал напівсферичного заземлювача спадає обернено пропорційно відстані від його центру (як у точкового заряду):
\begin{equation*}
	\phi(r) = \frac{\phi_0 D}{2}\cdot\frac1r.
\end{equation*}
Крокова напруга --- це різниця потенціалів між точками, де стоять дві ноги людини, на відстанях $r$ і $r+\ell$ від центру:
\begin{equation*}
	U_\text{крокова} = \phi(r) - \phi(r+\ell) \approx -\phi'(r)\,\ell = \frac{\phi_0 D}{2}\cdot\frac{\ell}{r^2} \approx 246\ \text{В}.
\end{equation*}

\begin{Attention}
	\emphx{Крокова напруга} --- важливе поняття техніки безпеки: людина, яка опинилась поблизу місця замикання проводу високої напруги на землю, не
	торкаючись самого проводу, все одно може отримати небезпечний удар струмом просто через різницю потенціалів між ступнями під час ходьби.
	Оскільки $U_\text{крокова}\propto 1/r^2$ спадає значно швидше, ніж сам потенціал $\phi\propto 1/r$, найнебезпечніша зона --- це ділянка \emphx{біля
		самого заземлювача}, а рекомендована поведінка --- віддалятися від місця замикання дрібними кроками (або стрибками, тримаючи ноги разом), щоб
	звести різницю потенціалів між ступнями до мінімуму.
\end{Attention}

%% --------------------------------------------------------
\section{Перехідні процеси в колі з конденсатором}
%% --------------------------------------------------------

Усе, що обговорювалося дотепер, стосувалося \emphx{стаціонарних} струмів, які не змінюються з часом. \emphx{Перехідними процесами} називають
процеси переходу від одного стаціонарного (чи взагалі відсутнього) режиму до іншого; найпростіший і найважливіший приклад --- заряджання та
розряджання конденсатора.

\begin{Attention}
	Закон Ома, отриманий для \emphx{постійних} струмів, залишається справедливим і для струмів, що \emphx{повільно} змінюються з часом --- за
	умови, що зміна відбувається не надто швидко (точніше --- повільно порівняно з часом встановлення електромагнітного поля в колі, тобто з часом
	пробігу електромагнітного сигналу вздовж кола). У такому разі миттєве значення струму практично однакове в усіх точках послідовно з'єднаної
	ділянки кола, і такі струми називають \emphx{квазістаціонарними}: до них можна застосовувати всі співвідношення для постійного струму, якщо
	тільки трактувати всі величини як їхні \emphx{миттєві} значення.
\end{Attention}

\subsection*{Розрядка конденсатора}

Розглянемо коло з конденсатора $C$, попередньо зарядженого до заряду $q_0$, резистора $R$ і ключа $S$, з'єднаних послідовно
(рис.~\ref{pic:RCDischarge}). Після замикання ключа, за законом Ома для кола, $RI = U$, де $U = q/C$ --- напруга на конденсаторі, а
$I = -dq/dt$ (струм тече в напрямку \emphx{зменшення} заряду конденсатора). Отримуємо диференціальне рівняння
\begin{equation}\label{eq:RCDischarge}
	\frac{dq}{dt} + \frac{q}{RC} = 0,
\end{equation}
розв'язком якого (з початковою умовою $q(0)=q_0$) є експоненційне спадання заряду:
\begin{equation}\label{eq:RCDischargeSolution}
	\tcbhighmath{
		q = q_0\, e^{-t/RC} = q_0\, e^{-t/\tau},
	}
\end{equation}
де величину $\tau = RC$ називають \emphx{часом релаксації} --- часом, за який заряд (а разом з ним струм і напруга) зменшується в $e\approx 2.72$
раза.

\begin{figure}[h!]\centering
	\localinput{rc-discharge.tikz}
	\caption{Розрядка конденсатора через резистор: коло та графік $q(t)$\label{pic:RCDischarge}}
\end{figure}

\subsection*{Зарядка конденсатора}

Аналогічно розглянемо зарядку конденсатора від джерела з ЕРС $\mathcal E$ через резистор $R$ (рис.~\ref{pic:RCCharge}). Закон Ома для кола тепер
$RI + U = \mathcal E$, причому $I = +dq/dt$ (струм тече в напрямку \emphx{зростання} заряду), звідки
\begin{equation}\label{eq:RCCharge}
	\frac{dq}{dt} + \frac{q}{RC} = \frac{\mathcal E}{R}.
\end{equation}
З початковою умовою $q(0) = 0$ розв'язок цього рівняння --- заряд, що монотонно й асимптотично наближається до граничного значення:
\begin{equation}\label{eq:RCChargeSolution}
	\tcbhighmath{
		q = q_{\max}\left(1 - e^{-t/\tau}\right), \qquad q_{\max} = \mathcal E C,
	}
\end{equation}
причому $q_{\max} = \mathcal E C$ --- це якраз рівноважний заряд конденсатора при $t\to\infty$ (коли струм у колі припиняється, а вся ЕРС падає на
конденсаторі), а стала часу $\tau = RC$ --- та сама, що й при розрядці.

\begin{figure}[h!]\centering
	\localinput{rc-charge.tikz}
	\caption{Зарядка конденсатора від джерела ЕРС через резистор: коло та графік $q(t)$\label{pic:RCCharge}}
\end{figure}

\begin{Attention}
	Стала часу $\tau = RC$ --- фундаментальна часова характеристика будь-якого $RC$-кола, яка визначає, наскільки швидко (чи повільно) коло реагує
	на зміну зовнішніх умов, і не залежить від того, заряджається конденсатор чи розряджається. Цікаво поєднати цей результат з опором середовища
	з формули~\eqref{eq:GroundResistanceGeneral}: якщо заряджену провідну кулю (власна ємність $C$) занурити в необмежене слабко провідне
	середовище (діелектрик з невеликою домішкою провідності $\lambda$) і дати їй "розрядитися" через це середовище, стала часу такого процесу
	дорівнює
	\begin{equation*}
		\tau = RC = \frac{\epsilon}{4\pi\lambda C}\cdot C = \frac{\epsilon}{4\pi\lambda},
	\end{equation*}
	тобто \emphx{не залежить від розміру й форми кулі}, а визначається \emphx{лише властивостями самого середовища} --- ще один прояв тієї самої
	аналогії між ємністю і опором, яку ми обговорювали вище.
\end{Attention}

%% --------------------------------------------------------
\section{Підсумки розділу}
%% --------------------------------------------------------

\subsection*{Основні поняття}
\begin{itemize}
	\item Сила струму: $\displaystyle I = \frac{dq}{dt}$; вектор густини струму: $\displaystyle j = \frac{dI}{dS_\perp}$; струм як потік вектора
	      $\vect j$: $\displaystyle I = \iint\limits_S \vect j\cdot d\vect S$.
	\item Зв'язок густини струму й густини заряду: $\vect j = \rho\vect u$.
\end{itemize}

\subsection*{Закони постійного струму}
\begin{itemize}
	\item Закон збереження заряду (рівняння неперервності): $\displaystyle \divg\vect j = -\frac{\partial\rho}{\partial t}$; для стаціонарного
	      струму: $\divg\vect j = 0$.
	\item Закон Ома в інтегральній формі: $\displaystyle I = \frac{U}{R}$, $\displaystyle R = \rho\frac{\ell}{S}$.
	\item Закон Ома в диференціальній формі: $\vect j = \lambda\Efield$; провідність металів за теорією Друде: $\displaystyle \lambda =
	      \frac{ne^2\tau}{2m_e}$.
	\item ЕРС: $\displaystyle \mathcal E = \frac{A^*}{q} = \int\limits_{(1)}^{(2)} \vect E^*\cdot d\vect r$; узагальнений закон Ома: $\vect j =
	      \lambda(\Efield+\vect E^*)$.
	\item Закон Ома для неоднорідної ділянки: $IR = \phi_1-\phi_2+\mathcal E$; для замкненого кола: $IR = \mathcal E$.
	\item Правила Кірхгофа: $\displaystyle \sum_k I_k = 0$ (для вузла); $\displaystyle \sum_m I_mR_m = \sum_n\mathcal E_n$ (для контуру).
	\item Закон Джоуля--Ленца: $w = \lambda E^2$ (диференціальна форма), $\mathcal P = I^2R$ (інтегральна форма).
\end{itemize}

\subsection*{Спеціальні задачі}
\begin{itemize}
	\item Опір необмеженого середовища між електродами через їхню взаємну ємність: $\displaystyle R = \frac{\epsilon}{4\pi\lambda C}$.
	\item $RC$-коло: $q(t) = q_0e^{-t/\tau}$ (розрядка) або $q_{\max}(1-e^{-t/\tau})$ (зарядка), $\tau = RC$.
\end{itemize}

\begin{Attention}
	Постійний струм --- це не просто "заряди, що рухаються", а \emphx{самоузгоджена} картина: розподіл поверхневих зарядів на провідниках створює
	поле $\Efield$ всередині них; це поле, за законом Ома, підтримує потрібну густину струму $\vect j$; а сам струм, через рівняння неперервності,
	вимагає, щоб розподіл зарядів (а отже, і поле) залишався незмінним у часі. Розірвати це самоузгодження й підтримати різницю потенціалів, попри
	неминучі теплові втрати (закон Джоуля--Ленца), можуть лише сторонні сили --- джерела ЕРС, без яких жоден постійний струм не міг би текти
	довше, ніж час розрядки конденсатора, утвореного самими провідниками кола.
\end{Attention}
